% ===== 原文 pp.188–202（PDF p.203–217）· 第 5 章 §5.3–5.4 · tr-06 =====
% [原文 p.188]
\begin{figure}[htbp]
\centering
\includegraphics[width=0.52\textwidth]{f5_6.png}
\figtitle{6}{}
\end{figure}

$\widetilde{\cM}_{\Sigma,k}\times\C^{\times}$，也就是 $\widetilde{\cM}_{\Sigma}\times\C^{\times}$。事实上，这个局部系统与 §2.2.3 中在 $\PP^2$ 情形给出的那个局部系统一致。

由 §2.2.4 我们知道，局部系统 $\cR$ 的秩为 $3$，并有多值截面基 $\Xi_0,\Xi_1,\Xi_2$，使得它们在 $e^{W_0(Q)/\hbar}\Omega$ 上的积分由 (2.38) 给出。于是积分
\[
  \int_{\Xi_i}e^{W_k(Q)/\hbar}\Omega
\]
在形式意义下有意义：只需写出
\[
  \exp(W_k(Q)/\hbar)=\exp(W_0(Q)/\hbar)\exp\bigl((W_k(Q)-W_0(Q))/\hbar\bigr),
\]
并把 $\exp((W_k(Q)-W_0(Q))/\hbar)$ 展成有限幂级数，完全如 §2.2.3 所做的那样。

本章的主要定理（将在 §5.5 证明）现在可以叙述了：

% [原文 p.189]
\begin{theorem}
我们可以写
\begin{equation}
  \sum_{i=0}^{2}\alpha^i\int_{\Xi_i}e^{W_k(Q)/\hbar}\Omega
  =\hbar^{-3\alpha}\sum_{i=0}^{2}\varphi_i\cdot(\alpha\hbar)^i, \tag{5.5}
\end{equation}
其中
\[
  \varphi_i(y_0,y_1,u_1,\ldots,u_k,\hbar^{-1})
  =\delta_{0,i}+\sum_{j=1}^{\infty}\varphi_{i,j}(y_0,y_1,u_1,\ldots,u_k)\hbar^{-j},
\]
对 $0\le i\le 2$，且
\[
  \begin{aligned}
    \varphi_{0,1}&=y_0\\
    \varphi_{1,1}&=y_1\\
    \varphi_{2,1}&=y_2:=\sum_{i=1}^{k}u_i.
  \end{aligned}
\]
此外，
\[
  \varphi_i=J_i^{\mathrm{trop}}(y_0,y_1,y_2).
\]
\end{theorem}

% [Tutor Note: p.189 (5.5) 中 $\hbar^{-3\alpha}$ 是形式记号 $\exp(-3\alpha\log\hbar)$（见第 2 章 §2.2.3 的约定）；因 $\alpha\in H^2(\PP^2)$ 满足 $\alpha^3=0$，$\hbar^{-3\alpha}=1-3\alpha\log\hbar+\frac{9}{2}(\alpha\log\hbar)^2$，故右端只出现 $(\alpha\hbar)^i$（$i\le2$）三项。核对时注意左端是 $\alpha^i$ 乘积分，右端是 $\hbar^{-3\alpha}(\alpha\hbar)^i$，两边的 $\alpha$ 幂次与 $\hbar$ 幂次不要弄混。]

事实上，这个定理证明了带后代的热带曲线计数与镜像对称之间的等价性。为使之精确，考虑如下陈述：

\begin{statement}[带后代的热带曲线计数]
\[
  J_{\PP^2}=J^{\mathrm{trop}}_{\PP^2}.
\]
\end{statement}

\begin{corollary}
定理 2.44 与陈述 5.19 等价。特别地，由于定理 2.44 已知为真，陈述 5.19 为真。
\end{corollary}

\begin{proof}
我们采用第 2 章的记号，其中 $\widetilde\cM=\widetilde\cM^B$ 是 B-模型模空间，$\pi:\check\cX\to\widetilde\cM$ 如 §2.2.3 所给，并且
\[
  W=t_0+W_0+t_2W_0^2,
\]
是 $\check\cX$ 上的势函数，在 $\widetilde\cM$ 的每个点的邻域内是通用的。另一方面，对固定的 $k$，我们有 §5.1.2 中构造的 $\kappa:\check\cX_{\Sigma,k}\to\widetilde\cM_{\Sigma,k}$。我们把 $\widetilde\cM$ 与 $\widetilde\cM_{\Sigma,k}$ 换成这些空间在 $0$ 处的芽，并把 $\check\cX$ 与 $\check\cX_{\Sigma,k}$ 换成这些芽的原像。于是 $(\check\cX,W)$ 是 $(\pi^{-1}(0),W_0)$ 的通用展开，而后者与 $(\kappa^{-1}(0),W_k(Q)|_{\kappa^{-1}(0)})$ 一致。因此有图表
\[
  \begin{array}{ccc}
    \check\cX_{\Sigma,k} & \xrightarrow{\eta} & \check\cX\\[4pt]
    {\scriptstyle\kappa}\big\downarrow & & \big\downarrow{\scriptstyle\pi}\\[4pt]
    \widetilde\cM_{\Sigma,k} & \xrightarrow{\xi} & \widetilde\cM
  \end{array}
\]
使得 $\eta|_{\kappa^{-1}(0)}:\kappa^{-1}(0)\to\pi^{-1}(0)$ 是恒等映射，且 $W\circ\eta=W_k(Q)$。

现在 $\widetilde\cM\times\C^{\times}$ 与 $\widetilde\cM_{\Sigma,k}\times\C^{\times}$ 都带有局部系统 $\cR$，其在 $(u,\hbar)$ 处的茎分别由
\[
  H_2(\pi^{-1}(u),\operatorname{Re}(W_0/\hbar)\ll 0;\C)
\]
或
\[
  H_2(\kappa^{-1}(u),\operatorname{Re}(W_k(Q)/\hbar)\ll 0;\C)
\]
给出。

% [原文 p.190]
为区分这两种情形，把 $\widetilde\cM\times\C^{\times}$ 上的局部系统写作 $\cR'$，则显然有 $\xi^{-1}\cR'\cong\cR$。

在 $\widetilde\cM$ 上有 B-模型的半无限 Hodge 结构变分，由 $\cE',\nabla',(\cdot,\cdot)_{\cE'}$ 以及由 Euler 向量场 $E'$ 确定的 $\mathrm{Gr}'$ 给出。类似地，在 $\widetilde\cM_{\Sigma,k}$ 上我们可以用完全相同的步骤构造半无限 Hodge 结构变分：此时不带撇号。于是 $\cE$ 的定义使得 $\Gamma(U,\cE)$ 由 $\cR^{\vee}=(\cR\otimes\cO_{\widetilde\cM_{\Sigma,k}\times\C^{\times}})^{\vee}$ 在形如 $U\times\{\hbar\mid |\hbar|<\epsilon\}$ 的开集上的截面组成，这些截面对应于 $\kappa^{-1}(U)\times\{\hbar\mid |\hbar|<\epsilon\}$ 上的形式 $f\Omega$，其中 $f$ 全纯，且限制到 $\kappa^{-1}(u)\times\{\hbar\}$ 时是代数的。联络 $\nabla$ 由 Gauss--Manin 联络诱导，配对由 $H_2$ 上的配对给出。唯一定义略有不同的是分次算子：
\[
  \mathrm{Gr}(s)=\nabla^{\mathrm{GM}}_{\hbar\partial_\hbar+E}(s)-s,
\]
其中
\[
  E=y_0\partial_{y_0}+3\partial_{y_1}-\sum_{i=1}^{k}u_i\partial_{u_i}.
\]
由构造显然有 $\cE=\xi^*\cE'$，并且 $\nabla',(\cdot,\cdot)_{\cE'}$ 拉回为 $\nabla$ 与 $(\cdot,\cdot)_{\cE}$。我断言 $\mathrm{Gr}'$ 也拉回为 $\mathrm{Gr}$。为验证这一点，注意
\[
  E=\kappa_*\Bigl(y_0\partial_{y_0}+\sum_{i=0}^{2}x_i\partial_{x_i}-\sum_{i=1}^{k}u_i\partial_{u_i}\Bigr),
\]
因此把 $E$ 看作 $\check\cX_{\Sigma,k}$ 上的向量场，有 $E(W_k(Q))=W_k(Q)$。事实上 $E(y_0)=y_0$，且
\[
  \begin{aligned}
    E(\operatorname{Mono}(h))
    &=E\bigl(\operatorname{Mult}(h)z^{\Delta(h)}u_{I(h)}\bigr)\\
    &=\bigl(|\Delta(h)|-\#I(h)\bigr)\operatorname{Mult}(h)z^{\Delta(h)}u_{I(h)}\\
    &=\operatorname{Mono}(h),
  \end{aligned}
\]
因为 $W_k(Q)$ 中出现的每个 $h$ 的 Maslov 指标都是二。于是把 $E$ 看作生活在 $\check\cX_{\Sigma,k}$ 上，$\eta_*E$ 保持 $W$，因为 $(\eta_*E)(W)=E(W\circ\eta)$，从而 $\eta_*E=E'$，这是由于 $E'$ 是 $\check\cX$ 上保持 $W$ 的唯一向量场。因此 $\mathrm{Gr}$ 是 $\mathrm{Gr}'$ 的拉回。

% [Tutor Note: p.190 核对点：$E(\operatorname{Mono}(h))=\operatorname{Mono}(h)$ 与 $E(W_k(Q))=W_k(Q)$ 等价于“$W_k(Q)$ 中出现的每个 $h$ 的 Maslov 指标都是二”，因 $z^{\Delta(h)}$ 在 $E$ 下的权为 $|\Delta(h)|$、$u_{I(h)}$ 的权为 $-\#I(h)$，故 $\operatorname{Mono}(h)$ 的权为 $|\Delta(h)|-\#I(h)=\mathrm{MI}(h)/2=1$（用 $\mathrm{MI}(h)=2(|\Delta(h)|-d)$ 且 $d=\#I(h)$）。]

现在 $\widetilde\cM$ 上的 Frobenius 流形结构由 $\cH'$ 的分裂 $\cH'_- \oplus\cE'_0$ 诱导，其中与 $\widetilde\cM$ 相关的所有记号仍带撇号。由于 $\cE'\otimes_{\cO_{\widetilde\cM}\{\hbar\}}\cO_{\widetilde\cM}\{\hbar,\hbar^{-1}\}$ 的平坦截面拉回为 $\cE\otimes_{\cO_{\widetilde\cM_{\Sigma,k}}\{\hbar\}}\cO_{\widetilde\cM_{\Sigma,k}}\{\hbar,\hbar^{-1}\}$ 的平坦截面，我们看到 $\cH$ 与 $\cH'$ 典范同构。当然 $\cE_0$ 与 $\cE'_0$ 也同构，因此我们可以使用分裂 $\cH=\cH_-\oplus\cE_0$，其中 $\cH_-=\cH'_-$。

令 $\Omega'_0\in\cH'$ 是在 $0\in\widetilde\cM$ 处取值 $[\Omega]$ 的平坦截面；它与在 $0\in\widetilde\cM_{\Sigma,k}$ 处取值 $[\Omega]$ 的平坦截面 $\Omega_0\in\cH$ 一致。我们有映射
\[
  \tau':(\hbar\cH'_-/\cH'_-)\otimes_\C\cO_{\widetilde\cM}\{\hbar\}\to\cE'
  \quad\text{与}\quad
  \tau:(\hbar\cH_-/\cH_-)\otimes_\C\cO_{\widetilde\cM_{\Sigma,k}}\{\hbar\}\to\cE
\]
（由 (2.16) 给出）；这些映射与 $\xi$ 的拉回交换。由命题 2.43 (3) 的证明我们知道，$\tau'([\Omega'_0]\otimes1)$ 可由某个 $f'\Omega$ 表示，其中 $f'$ 在 $\check\cX\times\{\hbar\mid |\hbar|<\epsilon\}$ 上全纯、在 $\pi$ 的纤维上正则、$f'|_{\pi^{-1}(0)}\equiv 1$，并且
\[
  \sum_{i=0}^{2}\alpha^i\int_{\Xi_i}e^{W/\hbar}f'\Omega
  =\hbar^{-3\alpha}\sum_{i=0}^{2}\varphi'_i(\mathbf t,\hbar^{-1})(\alpha\hbar)^i
\]

% [原文 p.191]
使得
\[
  \varphi'_i(\mathbf t,\hbar^{-1})=\delta_{0,i}+\sum_{j=1}^{\infty}\varphi'_{i,j}(\mathbf t)\hbar^{-j}.
\]
此外，这两个条件唯一确定由 $f'\Omega$ 表示的 $\cE'$ 的截面。对 $\tau([\Omega_0]\otimes1)$ 有同样结论：它可由某个 $f\Omega$ 表示，其中 $f$ 在 $\check\cX_{\Sigma,k}\times\{\hbar\mid|\hbar|<\epsilon\}$ 上全纯、在 $\kappa$ 的纤维上正则、$f|_{\kappa^{-1}(0)}\equiv1$，并且
\[
  \sum_{i=0}^{2}\alpha^i\int_{\Xi_i}e^{W_k(Q)/\hbar}f\Omega
  =\hbar^{-3\alpha}\sum_{i=0}^{2}\varphi_i(y_0,y_1,u_1,\ldots,u_k,\hbar^{-1})(\alpha\hbar)^i
\]
使得
\[
  \varphi=\delta_{0,i}+\sum_{j=1}^{\infty}\varphi_{i,j}(y_0,y_1,u_1,\ldots,u_k)\hbar^{-j}.
\]
现在要点是，定理 5.18 表明我们可以在 $\check\cX_{\Sigma,k}\times\C^{\times}$ 上取 $f\equiv1$。此外，由于 $[\Omega_0]=\xi^*[\Omega'_0]$，有 $\tau([\Omega_0]\otimes1)=\xi^*\tau'([\Omega'_0]\otimes1)$。于是若 $f'\Omega$ 如上表示 $\tau'([\Omega'_0]\otimes1)$，我们看到
\[
  \varphi_i=\varphi'_i\circ\xi.
\]
特别地，若 $y'_0,y'_1,y'_2$ 是 $\widetilde\cM$ 上的平坦坐标，即 $y'_i=\varphi'_{i,1}$，则
\[
  y_i=y'_i\circ\xi.
\]

现在假设定理 2.44 为真。则特别地，命题 2.45 的结论成立，且 $\varphi'_i=J_i$，故 $\varphi_i=J_i\circ\xi$。但由定理 5.18，$\varphi_i=J_i^{\mathrm{trop}}$。于是等同坐标 $y_i$ 与 $y'_i$，我们看到 $J_i=J_i^{\mathrm{trop}}$，即陈述 5.19 成立。

反之，假设陈述 5.19 成立。我们只需证明命题 2.45 的结论成立。注意我们已知 (1)：它是 $\widetilde\cM$ 上 Frobenius 流形结构构造的一部分。然后陈述 5.19 给出 (2)。由于
\[
  E_A=y_0\partial_{y_0}+3\partial_{y_1}-y_2\partial_{y_2},
\]
于是当 $y_2=u_1+\cdots+u_k$ 时，
\[
  E=y_0\partial_{y_0}+3\partial_{y_1}-\sum_{i=1}^{k}u_i\partial_{u_i}
  =y_0\partial_{y_0}+3\partial_{y_1}-y_2\partial_{y_2}
\]
与 $E_A$ 一致。这给出 (3)。最后对 (4)，我们在命题 2.43 (2) 的证明中论述过 $\bigl((-) ^*(\hbar^{-3\alpha}\alpha^j),\hbar^{-3\alpha}\alpha^k\bigr)$ 是单值的，而由于对 $\hbar$ 的依赖是 $\log\hbar$ 的函数，它必须与 $\hbar$ 无关。特别地，利用 $\hbar^{-3\alpha}\alpha^j=(1-3\alpha\log\hbar+(9/2)(\alpha\log\hbar)^2)\alpha^j$，我们看到
\[
  \bigl((-) ^*\hbar^{-3\alpha}\alpha^0,\hbar^{-3\alpha}\alpha^2\bigr)
  =\bigl((-) ^*\alpha^0,\alpha^2\bigr).
\]
回忆 $1,\alpha,\alpha^2$ 是 $\Xi_0,\Xi_1,\Xi_2$ 的对偶基。由注 2.42 中 $\Xi_i$ 的具体形式以及例 2.34 中的相交数，我们有
\[
  ((-)^*\Xi_0,\Xi_2)=-\frac{1}{(2\pi\ii)^2},
\]
故
\[
  ((-)^*\alpha^0,\alpha^2)=-(2\pi\ii)^2,
\]
从而
\[
  (\hbar^{-3\alpha}\alpha^0,\hbar^{-3\alpha}\alpha^2)_{\cE}=\hbar^{-2}.
\]

% [原文 p.192]
把 (2.39) 用两次并利用 (2.40)，我们看到
\[
  (\hbar^{-3\alpha}\alpha^i,\hbar^{-3\alpha}\alpha^j)_{\cE}=
  \begin{cases}
    0 & i+j>2,\\
    (-1)^i\hbar^{-2} & i+j=2,\\
    (\text{常数})\hbar^{-2} & i+j<2.
  \end{cases}
\]
于是
\[
  \begin{aligned}
    (\hbar^{-3\alpha}(\hbar\alpha)^i,\hbar^{-3\alpha}(\hbar\alpha)^j)_{\hbar\cH_-/\cH_-}
    &=\bigl((-\hbar)^i\hbar^j(\hbar^{-3\alpha}\alpha^i,\hbar^{-3\alpha}\alpha^j)_{\cE}\bigr)\big|_{\hbar=\infty}\\
    &=\delta_{i,2-j}\\
    &=\int_{\PP^2}T_i\cup T_j.
  \end{aligned}
\]
这给出命题 2.45 的 (4)，从而给出定理 2.44。
\end{proof}

% [Tutor Note: p.192 核对锚点：$i+j>2$ 时为零只用 $\alpha^3=0$；而 $(\hbar^{-3\alpha}(\hbar\alpha)^i,\hbar^{-3\alpha}(\hbar\alpha)^j)_{\hbar\cH_-/\cH_-}$ 取 $\hbar\to\infty$ 的极限后等于 $\delta_{i,2-j}$，正是 $\int_{\PP^2}T_i\cup T_j$ 的值——注意指标 $i,j$ 在 $(2.39)$ 的对偶配对下是反转的，$T_i$ 与 $\alpha^i$ 的对偶关系为 $(T_i,\alpha^j)=\delta_{ij}$。]

% [Tutor Note: p.191 原文 “$\varphi=\delta_{0,i}+\sum_j\varphi_{i,j}\hbar^{-j}$” 的左端漏了下标 $i$（应为 $\varphi_i$）：依据是本页上一条显示式左端为 $\varphi_i(\ldots)$、正文 Theorem 5.18 的写法，以及右端 $\delta_{0,i}$ 与 $\varphi_{i,j}$ 都带 $i$。此为原文笔误，照译以存真。]

定理 5.15 与 5.18 的证明将在 §5.5 给出。这些证明需要显式计算积分 $\int_{\Xi_i}e^{W_k(Q)/\hbar}\Omega$。虽然一般说来得到这些积分的显式答案并不困难，但从这些答案中提取有用的组合学结果却相当困难。不过对 $\Xi_0$ 的计算相当容易，要看它被显式地做出来，读者可参看 [42] 的命题 3.13。然而即使对 $\Xi_0$ 的论证也将被包含在更一般的计算之中，所以我们不在此重复。

\subsection{形变 $Q$ 与 $P_1,\ldots,P_k$}

本章的最终目标是计算积分 $\int_{\Xi_i}e^{W_k(Q)/\hbar}\Omega$。这些积分既涉及 $Q$ 的显式选取，也涉及 $P_1,\ldots,P_k$ 的隐含选取。然而我们期望这些积分与这些选取无关。我们需要证明这一点，而为证明这一点，我们需要 $W_k(Q)$ 随 $Q$ 或 $P_1,\ldots,P_k$ 变化而变化的显式公式。事实上，我们将看到存在一个“跨墙公式”，它说 $W_k(Q)$ 的变化由一个自同构给出，该自同构作用在 $W_k(Q)$ 所生活的环上。这些跨墙发生于某些圆盘模空间维数不正确之时。这些跨墙似乎对该理论至关重要，并将在下一章以不同的面貌再次出现。

为引出讨论，我们先非正式地探讨 Maslov 指标零圆盘所起的作用。

\subsubsection{穿越 Maslov 指标零圆盘}

在本讨论中，我们固定一般位置的 $P_1,\ldots,P_k$ 并让 $Q$ 变化。由引理 5.6，以一般位置点 $Q$ 为边界的 Maslov 指标零圆盘不存在，因为这类圆盘的模空间维数为 $-1$。另一方面，不带指定边界的 Maslov 指标零圆盘的模空间是一维的，所以我们期望这类圆盘存在。

注意这个模空间的一维性来自如下事实。若有一个以某点 $Q$ 为边界的 Maslov 指标零圆盘 $g:\Gamma'\to M_\R$，带出射边 $E_{\mathrm{out}}$，我们只需把 $Q$ 沿平行于 $g(E_{\mathrm{out}})$ 的方向移动，就可使 $g(E_{\mathrm{out}})$ 变短或变长。因此，通过把 $g(E_{\mathrm{out}})$ 延伸为无界，来记录由此得到的整族一参数 Maslov 指标零圆盘是有意义的；见图 7。左图显示一个 Maslov 指标零热带圆盘，而在右图中可沿标为 $g(E_{\mathrm{out}})$ 的射线把 $Q$ 在一参数族中变动。这里 $g:\Gamma\to M_\R$ 是由左边热带圆盘出发、把圆盘出射边无限延伸而得到的热带曲线。

% [原文 p.193]
我们仍把这条无界边称作 $E_{\mathrm{out}}$。注意虽然这是热带曲线，但它不是 $\PP^2$ 中的曲线，因为 $g(E_{\mathrm{out}})$ 未必平行于 $\PP^2$ 的扇 $\Sigma$ 中的任何射线。然而任何 $Q\in g(E_{\mathrm{out}})$ 都是某个 Maslov 指标零圆盘的边界。我们称热带圆盘的这类延伸为热带树：这一点将在 §5.4.3 精确化。

\begin{figure}[htbp]
\centering
\includegraphics[width=0.86\textwidth]{f5_7.png}
\figtitle{7}{}
\end{figure}

现在假设我们有这样一棵热带树 $g$，它给出单参数族的 Maslov 指标零圆盘。考虑 $g(E_{\mathrm{out}})$ 两侧的两个点 $Q_1,Q_2$。假设有一个以 $Q_1$ 为边界的 Maslov 指标二圆盘。当我们从 $Q_1$ 移到 $Q_2$ 时，这个圆盘可能发生什么？

我们在此在直观层面描述各种可能性；精确结果将在 §5.4.5 给出。

从以 $Q_1$ 为边界、出射边为 $E_{\mathrm{out},h}$ 的 Maslov 指标二圆盘 $h$ 开始。令 $m\in M$ 是 $h(E_{\mathrm{out},h})$ 上指向远离 $Q_1$ 方向的本原切向量。有三种可能：

\textit{情形 1.} 切向量 $m$ 可能指向远离 $g(E_{\mathrm{out}})$ 的方向。此时我们可尝试形变该 Maslov 指标二圆盘，使它的边界从 $Q_1$ 连续移动到 $Q_2$。一般地（我们将看到）这是可以做到的。这给出一个以 $Q_2$ 为边界的 Maslov 指标二圆盘。但还有一件事可做：把圆盘形变到以 $h(E_{\mathrm{out}})$ 上的点 $Q'$ 为边界之后，我们可以取以 $Q'$ 为边界的 Maslov 指标零圆盘，把它与以 $Q'$ 为边界的 Maslov 指标二圆盘粘接，从而得到一种新的 Maslov 指标二圆盘。然后按平衡条件确定权与方向，再加上一条新的出射边；见图 8。于是单个以 $Q_1$ 为边界的 Maslov 指标二圆盘给出两个以 $Q_2$ 为边界的 Maslov 指标二圆盘：这正是我们在此要研究的跨墙现象的基本特征。

令 $h_1,h_2$ 是这样构造出的两个以 $Q_2$ 为边界的 Maslov 指标二圆盘，其中 $h_2$ 是由粘接过程得到的。则显然 $\operatorname{Mono}(h_1)=\operatorname{Mono}(h)$，但 $h_2$ 比来自 $h$ 与 $g$ 的顶点多出一个顶点 $V$，其重数为
\[
  w(E_{\mathrm{out}})w(E_{\mathrm{out},h})|m\wedge m'|,
\]
其中 $m'$ 是 $g(E_{\mathrm{out}})$ 的本原切向量。注意用圆盘 $h$ 的次数 $\Delta(h)$ 表示，我们实际上有 $w(E_{\mathrm{out},h})m=r(\Delta(h))$，这只需把每个顶点处的平衡条件求和即可看出。类似地，若把 $\Delta(g)$ 记作给出 $g$ 的 Maslov 指标零圆盘次数，则可写

% [原文 p.194]
\begin{figure}[htbp]
\centering
\includegraphics[width=0.92\textwidth]{f5_8.png}
\figtitle{8}{}
\end{figure}

$w(E_{\mathrm{out}})m'=r(\Delta(g))$。于是
\begin{equation}
  \begin{aligned}
    \operatorname{Mono}(h_1)+\operatorname{Mono}(h_2)
    &=\operatorname{Mono}(h)\bigl(1+|r(\Delta(g))\wedge r(\Delta(h))|\operatorname{Mono}(g)\bigr)\\
    &=\operatorname{Mono}(h)(1+\operatorname{Mono}(g))^{|r(\Delta(g))\wedge r(\Delta(h))|}.
  \end{aligned}
  \tag{5.6}
\end{equation}
第二个等号来自 $u_{I(g)}^2=0$。

\textit{情形 2.} 第二种情形是 $m$ 指向 $g(E_{\mathrm{out}})$。此时我们可以把 $Q_1$ 与 $Q_2$ 的角色互换，于是用同样的论证，$W_k(Q_1)$ 包含来自 $W_k(Q_2)$ 中项 $\operatorname{Mono}(h)$ 的两个单项式 $\operatorname{Mono}(h_1)+\operatorname{Mono}(h_2)$。此时可将 (5.6) 改写为
\begin{equation}
  \operatorname{Mono}(h)=(\operatorname{Mono}(h_1)+\operatorname{Mono}(h_2))(1+\operatorname{Mono}(g))^{-|r(\Delta(g))\wedge r(\Delta(h))|}.
  \tag{5.7}
\end{equation}

\textit{情形 3.} 若 $m$ 平行于 $g(E_{\mathrm{out}})$，则预期在穿过 $g(E_{\mathrm{out}})$ 时不可能粘上一份 Maslov 指标零圆盘，于是当 $Q$ 穿过 $g(E_{\mathrm{out}})$ 时，来自这样的 $h$ 的贡献不变。

我们将于 §5.4.5 严格证明上述讨论是准确的；先验地看，Maslov 指标二圆盘能否随意形变并不清楚。但让我们考虑公式 (5.6) 与 (5.7) 的推论。定义 $N=\operatorname{Hom}_\Z(M,\Z)$ 中的元素 $n_g$

% [Tutor Note: p.194 (5.6) 的第二个等号用到 $u_{I(g)}^2=0$：令 $N=|r(\Delta(g))\wedge r(\Delta(h))|\ge1$，则 $(1+\operatorname{Mono}(g))^N$ 中凡含 $\operatorname{Mono}(g)^{\ge2}$ 的项都含因子 $u_{I(g)}^{\ge2}=0$，故只剩 $1+N\operatorname{Mono}(g)$，即得 (5.6)。而 (5.7) 是同一式在“反向穿越”下的逆，指数取负。核对 $w(E_{\mathrm{out}})m'=r(\Delta(g))$ 与 $w(E_{\mathrm{out},h})m=r(\Delta(h))$：它们分别由 $g$ 与 $h$ 各顶点平衡条件求和得到。]

% [原文 p.195]
为
\[
  m\mapsto
  \begin{cases}
    |r(\Delta(g))\wedge m| & \text{若 } m \text{ 指向远离 } g(E_{\mathrm{out}}) \text{ 的方向},\\
    -|r(\Delta(g))\wedge m| & \text{若 } m \text{ 指向 } g(E_{\mathrm{out}}),\\
    0 & \text{若 } m \text{ 平行于 } g(E_{\mathrm{out}}).
  \end{cases}
\]
容易看出这是线性的：事实上，它可以写成复合
\[
  M\to\bigwedge{}^2M\xrightarrow{\ \cong\ }\Z,
\]
其中第一个映射是
\[
  m\mapsto r(\Delta(g))\wedge m,
\]
而第二个是两个同构之一的选取，使得当 $m$ 指向远离 $g(E_{\mathrm{out}})$ 时 $r(\Delta(g))\wedge m$ 被映到正整数。

于是我们定义环 $\C[T_\Sigma]\otimes_\C R_k[[y_0]]$ 的自同构 $\theta_g$ 为
\[
  \theta_g(z^m)=z^m(1+\operatorname{Mono}(g))^{\langle n_g,r(m)\rangle}
\]
对 $m\in T_\Sigma$，以及
\[
  \theta_g(a)=a
\]
对 $a\in R_k[[y_0]]$。注意 $\theta_g$ 的逆由
\[
  \theta_g^{-1}(z^m)=z^m(1+\operatorname{Mono}(g))^{-\langle n_g,r(m)\rangle}
\]
给出，因此它是自同构。然后把 (5.6) 与 (5.7) 以及第三种情形的讨论合起来，得到
\begin{equation}
  W_k(Q_2)=\theta_g(W_k(Q_1)). \tag{5.8}
\end{equation}
我们尚未给出这一点的严格证明，甚至尚未给出严格的陈述。我们将在定理 5.35 中完整地做到。现在要领会的主要之点是：Maslov 指标零圆盘对 $W_k(Q)$ 随 $Q$ 变化负责，而 $W_k(Q)$ 的变化方式受到很好的控制，可用由 Maslov 指标零圆盘确定的 $\C[T_\Sigma]\otimes_\C R_k[[y_0]]$ 的自同构来表达。

当点 $P_1,\ldots,P_k$ 移动时，我们将发现类似但更复杂的行为。

现在我们更严格地探讨这些想法，首先描述在跨墙中出现的自同构群。

\subsubsection{自同构}

现在的目标是指定 $\C[T_\Sigma]\otimes_\C R_k[[y_0]]$ 的自同构群的一个子群。这些自同构特别地将是辛同胚，即它们保持 (5.3) 中定义的相对全纯二形式 $\Omega$。然而我们不考虑所有辛同胚，而考虑具有特殊性质的一个子群，它是 Kontsevich 与 Soibelman 在 [70] 中首先定义的一个群的变体。这些性质要到下一章才会被充分探讨：此处的讨论可看作一个热身。

在本小节中，$M\cong\Z^2$，$\Sigma$ 是 $M_\R$ 中定义非异环面曲面的任意完备扇。我们有标准正合列
\[
  0\to K_\Sigma\to T_\Sigma\xrightarrow{\ r\ }M\to 0.
\]
首先，我们定义 $\C[T_\Sigma]\otimes_\C R_k[[y_0]]$ 的对数导子模为模
\[
  \Theta(\C[T_\Sigma]\otimes_\C R_k[[y_0]])
  :=\operatorname{Hom}_\Z(M,\C[T_\Sigma]\otimes_\C R_k[[y_0]])
  =(\C[T_\Sigma]\otimes_\C R_k[[y_0]])\otimes_\Z N.
\]

% [原文 p.196]
元素 $f\otimes n$ 写作 $f\partial_n$，它作为 $\C[T_\Sigma]\otimes_\C R_k[[y_0]]$ 上的导子（在 $\C[K_\Sigma]\otimes_\C R_k[[y_0]]$ 之上）作用，通过
\[
  f\partial_n(z^m)=f\langle n,r(m)\rangle z^m.
\]
给定 $\xi\in\mathfrak m_{R_k}\Theta(\C[T_\Sigma]\otimes_\C R_k[[y_0]])$，其中 $\mathfrak m_{R_k}=(u_1,\ldots,u_k)$ 是 $R_k$ 的极大理想，我们定义
\[
  \exp(\xi)\in\operatorname{Aut}(\C[T_\Sigma]\otimes_\C R_k[[y_0]])
\]
为
\[
  \exp(\xi)(a)=a+\sum_{i=1}^{\infty}\frac{\xi^i(a)}{i!}.
\]
由对 $\xi$ 的假设，这是有限和，因为 $\mathfrak m_{R_k}^{k+1}=0$。容易验证 $\exp(\xi)$ 确实是环自同构。

% [Tutor Note: p.196 $\exp(\xi)$ 为有限和依赖于 $\mathfrak m_{R_k}^{k+1}=0$，即 $R_k$ 是加厚环（$u_i$ 幂零，$u_I=\prod_{i\in I}u_i$，$|I|\le k$）；这也解释了 (5.9) 中只出现一次 $u_I$——更高次幂全为零。核对 (5.9) 时注意 $\partial_n$ 作用在 $z^{m'}$ 上给出 $\langle n,r(m')\rangle z^{m'}$，故括号内是 $\langle n,r(m')\rangle$ 而非 $\langle n,m'\rangle$。]

对一般的对数导子，描述这个自同构可能很困难。不过这里有一个简单的例子：取 $m\in T_\Sigma$ 且 $r(m)\ne0$，$n\in N$ 满足 $\langle n,r(m)\rangle=0$，$I\subseteq\{1,\ldots,k\}$ 非空，$c\in\C$。则
\begin{equation}
  \exp(cu_Iz^m\partial_n)(z^{m'})=z^{m'}+cu_Iz^m\partial_nz^{m'}
  =z^{m'}(1+cu_I\langle n,r(m')\rangle z^m). \tag{5.9}
\end{equation}
这恰好是 §5.4.1 中出现的自同构 $\theta_g$ 的形式。

对数导子模上有自然的 Lie 括号：
\[
  \begin{aligned}
    [z^m\partial_n,z^{m'}\partial_{n'}]
    &:=z^m\partial_n(z^{m'})\partial_{n'}-z^{m'}\partial_{n'}(z^m)\partial_n\\
    &=z^{m+m'}\bigl(\langle n,r(m')\rangle\partial_{n'}-\langle n',r(m)\rangle\partial_n\bigr).
  \end{aligned}
\]
这就是通常的 Lie 括号。

现在令
\[
  \mathfrak v_{\Sigma,k}
  :=\bigoplus_{\substack{m\in T_\Sigma\\ r(m)\ne0}}z^m\bigl(\mathfrak m_{R_k}\otimes r(m)^{\perp}\bigr)
  \subseteq\Theta(\C[T_\Sigma]\otimes_\C R_k[[y_0]]).
\]
换言之，$\mathfrak v_{\Sigma,k}$ 由形如 $z^m\partial_n$（其中 $r(m)\ne0$ 且 $\langle n,r(m)\rangle=0$）的对数导子的 $\mathfrak m_{R_k}$-线性组合组成。

本讨论的关键结论是：$\mathfrak v_{\Sigma,k}$ 在 Lie 括号下封闭。事实上，给定满足上述两个要求的 $m,n$ 与 $m',n'$，我们注意
\[
  [z^m\partial_n,z^{m'}\partial_{n'}]
  =z^{m+m'}\partial_{\langle n,r(m')\rangle n'-\langle n',r(m)\rangle n}.
\]
现在
\[
  \begin{aligned}
    \bigl\langle\langle n,r(m')\rangle n'-\langle n',r(m)\rangle n,\ r(m+m')\bigr\rangle
    &=\langle n,r(m')\rangle\langle n',r(m)\rangle-\langle n',r(m)\rangle\langle n,r(m')\rangle\\
    &=0,
  \end{aligned}
\]
因为按假设 $\langle n,r(m)\rangle=\langle n',r(m')\rangle=0$。另一方面，在 $r(m+m')=0$ 的情形，$r(m)=-r(m')$ 且 $\langle n,r(m')\rangle=\langle n',r(m)\rangle=0$。于是在此情形 $[z^m\partial_n,z^{m'}\partial_{n'}]=0$ 无论如何成立。这表明 $\mathfrak v_{\Sigma,k}$ 是对数导子模的 Lie 子代数。

现在令
\[
  \mathbb V_{\Sigma,k}=\{\exp(\xi)\mid\xi\in\mathfrak v_{\Sigma,k}\}.
\]
由于 $\mathfrak v_{\Sigma,k}$ 在 Lie 括号下封闭（即 $\mathfrak v_{\Sigma,k}$ 是 Lie 代数）这一观察，$\mathbb V_{\Sigma,k}$ 是一个群，其乘法由 Baker--Campbell--Hausdorff 公式给出。注意由于 $\mathfrak v_{\Sigma,k}$ 由导子 $cu_Iz^m\partial_n$

% [原文 p.197]
（其中 $r(m)\ne0$，$\langle n,r(m)\rangle=0$）生成，$\mathbb V_{\Sigma,k}$ 由形如 $\exp(cu_Iz^m\partial_n)$ 的自同构生成，如 (5.9) 所述。

不难验证 $\exp(cu_Iz^m\partial_n)$ 保持 $\Omega$；这虽可直接验证，但很快也会自然得到。事实上，这个群在 [70] 中最初的版本被定义为 Hamilton 辛同胚群。

这个描述如下。我们固定相对全纯二形式 $\Omega$，如 (5.3)。它由 $\bigwedge^2M$ 的一个生成元的选取（即一个等同 $\bigwedge^2M\cong\Z$）确定。给定元素 $f\in\C[T_\Sigma]\otimes_\C R_k[[y_0]]$，我们记 $X_f$ 为关于辛形式 $\Omega$ 的相应 Hamilton 向量场。换言之，这是满足 $\iota(X_f)\Omega=df$ 的向量场 $X_f$。这里我们考虑的是关于映射 $\kappa$ 的相对微分。所以 $df$ 是对数导子模的对偶中的元素，即 $(\C[T_\Sigma]\otimes_\C R_k[[y_0]])\otimes_\Z M$ 的元素。我们把 $f\otimes m$ 写作 $f\,\dd\log m$，且 $\dd(z^m)=z^m\dd\log r(m)$。

为描述 $X_{z^m}$，我们引入由 $\Omega$ 所给出的等同 $\bigwedge^2M\cong\Z$ 诱导的 $N$ 与 $M$ 的一个等同。元素 $m\in M$ 诱导元素 $X_m\in N$，它由复合
\[
  \begin{array}{c}
    X_m:M\to\bigwedge{}^2M\xrightarrow{\ \cong\ }\Z\\
    m'\mapsto m\wedge m'
  \end{array}
\]
给出。于是若 $f=z^m$，则 $X_f$ 是向量场 $-z^mX_{r(m)}$。事实上，取 $e_1,e_2$ 为 $M$ 的有向基，写 $\Omega=\dd\log e_1\wedge\dd\log e_2$，且 $r(m)=a_1e_1+a_2e_2$，我们看到
\[
  \begin{aligned}
    -z^m\iota(X_{r(m)})\Omega
    &=-z^m\bigl(\langle X_{r(m)},e_1\rangle\dd\log e_2-\langle X_{r(m)},e_2\rangle\dd\log e_1\bigr)\\
    &=-z^m(-a_2\dd\log e_2-a_1\dd\log e_1)\\
    &=z^m\dd\log r(m)\\
    &=\dd(z^m).
  \end{aligned}
\]
注意 $\langle X_m,m\rangle=0$ 是平凡的。于是可以看出，事实上 $\mathfrak v_{\Sigma,k}$ 可描述为与形如 $\sum_ic_iu_{I_i}z^{m_i}$ 的函数 $f$ 相伴随的 Hamilton 向量场之集，其中 $c_i\in\C$，$I_i\subseteq\{1,\ldots,k\}$ 是非空指标集，$m_i\in T_\Sigma$ 且 $r(m_i)\ne0$。当然，Hamilton 向量场的指数映射保持 $\Omega$，因此 $\mathbb V_{\Sigma,k}$ 是辛同胚群。

为将来使用，我们有如下标准引理：

\begin{lemma}
若 $f\in\mathfrak m_{R_k}(\C[T_\Sigma]\otimes_\C R_k[[y_0]])$ 且 $\theta\in\mathbb V_{\Sigma,k}$，则
\[
  \theta\circ X_f\circ\theta^{-1}=X_{\theta(f)}.
\]
\end{lemma}

\begin{proof}
我们可以假设 $\theta$ 是 $\mathbb V_{\Sigma,k}$ 的生成元，即
\[
  \theta(z^{m''})=z^{m''}(1+r(m)\wedge r(m'')cu_Iz^m)
\]

% [原文 p.198]
对某个 $m\in T_\Sigma$，$r(m)\ne0$。由线性性，可假设 $f=z^{m'}$ 对某个 $m'$。则计算得
\[
  \begin{aligned}
    &(\theta\circ X_f\circ\theta^{-1})(z^{m''})\\
    &\quad=-\bigl(\theta\circ(z^{m'}X_{r(m')})\bigr)\bigl(z^{m''}-cu_Ir(m)\wedge r(m'')z^{m+m''}\bigr)\\
    &\quad=-\theta\bigl(r(m')\wedge r(m'')z^{m'+m''}\\
    &\qquad\qquad-cu_I[r(m')\wedge(r(m)+r(m''))][r(m)\wedge r(m'')]z^{m+m'+m''}\bigr)\\
    &\quad=-r(m')\wedge r(m'')z^{m'+m''}\\
    &\qquad-cu_I[r(m')\wedge r(m'')][r(m)\wedge(r(m')+r(m''))]z^{m+m'+m''}\\
    &\qquad+cu_I[r(m')\wedge(r(m)+r(m''))][r(m)\wedge r(m'')]z^{m+m'+m''}\\
    &\quad=-r(m')\wedge r(m'')z^{m'+m''}\\
    &\qquad-cu_I[r(m)\wedge r(m')][r(m+m')\wedge r(m'')]z^{m+m'+m''}.
  \end{aligned}
\]
另一方面，
\[
  \begin{aligned}
    X_{\theta(f)}(z^{m''})
    &=X_{z^{m'}+r(m)\wedge r(m')cu_Iz^{m+m'}}(z^{m''})\\
    &=-r(m')\wedge r(m'')z^{m'+m''}\\
    &\qquad-cu_I[r(m)\wedge r(m')][r(m+m')\wedge r(m'')]z^{m+m'+m''}.
  \end{aligned}
\]
\end{proof}

\subsubsection{散射图}

现在引入一种图示方法，我们称之为散射图方法，用来记录群 $\mathbb V_{\Sigma,k}$ 中的计算。照常固定 $M\cong\Z^2$，$\Sigma$ 是 $M_\R$ 中定义非异环面曲面的完备扇。

\begin{definition}
固定 $k\ge0$。
\begin{enumerate}[label=(\arabic*),leftmargin=*]
\item \textit{射线}或\textit{直线}是一个对 $(\mathfrak d,f_{\mathfrak d})$，使得
\begin{itemize}
\item $\mathfrak d\subseteq M_\R$ 由
\[
  \mathfrak d=m'_0-\R_{\ge0}r(m_0)
\]
（若 $\mathfrak d$ 是射线）或
\[
  \mathfrak d=m'_0-\R r(m_0)
\]
（若 $\mathfrak d$ 是直线）给出，其中 $m'_0\in M_\R$，$m_0\in T_\Sigma$ 且 $r(m_0)\ne0$。集合 $\mathfrak d$ 称为该直线或射线的\textit{支撑}。若 $\mathfrak d$ 是射线，则 $m'_0$ 称为该射线的\textit{初始点}，记作 $\operatorname{Init}(\mathfrak d)$。
\item $f_{\mathfrak d}\in\C[z^{m_0}]\otimes_\C R_k\subseteq\C[T_\Sigma]\otimes_\C R_k[[y_0]]$。
\item $f_{\mathfrak d}\equiv1\bmod(u_1,\ldots,u_k)z^{m_0}$。
\end{itemize}
\item \textit{散射图} $\mathfrak D$ 是直线与射线的有限集合。
\end{enumerate}
若 $\mathfrak D$ 是散射图，我们写
\[
  \operatorname{Supp}(\mathfrak D):=\bigcup_{\mathfrak d\in\mathfrak D}\mathfrak d\subseteq M_\R
\]
以及
\[
  \operatorname{Sing}(\mathfrak D):=\bigcup_{\mathfrak d\in\mathfrak D}\partial\mathfrak d
  \ \cup\bigcup_{\substack{\mathfrak d_1,\mathfrak d_2\\ \dim\mathfrak d_1\cap\mathfrak d_2=0}}\mathfrak d_1\cap\mathfrak d_2.
\]
\end{definition}

% [Tutor Note: p.198 记号说明：原文用花体 $\mathfrak d$ 表示射线/直线、$\mathfrak D$ 表示散射图，后文的 $\operatorname{Supp}$、$\operatorname{Sing}$、$\operatorname{Init}$ 均按此记号；$\mathfrak d=m'_0-\R_{\ge0}r(m_0)$ 表示以 $m'_0$ 为初始点、向内方向平行于 $-r(m_0)$ 的射线（直线则取整条 $\R r(m_0)$ 的平移）。$\partial\mathfrak d$ 是拓扑边界：射线时即初始点，直线时为空。]
% [原文 p.199]
这里 $\partial\mathfrak d=\{\operatorname{Init}(\mathfrak d)\}$ 若 $\mathfrak d$ 是射线，若 $\mathfrak d$ 是直线则为空。

\begin{figure}[htbp]
\centering
\includegraphics[width=0.86\textwidth]{f5_9.png}
\figtitle{9}{左图路径在 $M_\R\setminus\operatorname{Sing}(\mathfrak D)$ 中同伦于右图路径。自同构 $\theta_{\gamma,\mathfrak D}$ 在左右两图中相同，因为对同一直线的两次穿越相互抵消。}
\end{figure}

\begin{construction}
给定光滑浸入 $\gamma:[0,1]\to M_\R\setminus\operatorname{Sing}(\mathfrak D)$，端点不在 $\operatorname{Supp}(\mathfrak D)$ 中，且 $\gamma$ 横截地相交 $\mathfrak D$ 的元素，我们可以定义一个环自同构 $\theta_{\gamma,\mathfrak D}\in\mathbb V_{\Sigma,k}$，称为 $\mathfrak D$ 的 $\gamma$-有序乘积。确切地说，可以找到数
\[
  0<t_1\le t_2\le\cdots\le t_s<1
\]
与元素 $\mathfrak d_i\in\mathfrak D$，使得 $\gamma(t_i)\in\mathfrak d_i$，且当 $t_i=t_j$、$i\ne j$ 时 $\mathfrak d_i\ne\mathfrak d_j$，而 $s$ 取得尽可能大。然后对每个 $i$，定义 $\theta_{\gamma,\mathfrak d_i}\in\mathbb V_{\Sigma,k}$ 为
\[
  \begin{aligned}
    \theta_{\gamma,\mathfrak d_i}(z^m)&=z^mf_{\mathfrak d_i}^{\langle n_0,r(m)\rangle}\\
    \theta_{\gamma,\mathfrak d_i}(a)&=a
  \end{aligned}
\]
对 $m\in T_\Sigma$，$a\in R_k[[y_0]]$，其中 $n_0\in N$ 取为本原的、零化 $\mathfrak d_i$ 的切空间，并最终由符号约定
\[
  \langle n_0,\gamma'(t_i)\rangle<0
\]
完全确定。然后我们定义
\[
  \theta_{\gamma,\mathfrak D}=\theta_{\gamma,\mathfrak d_s}\circ\cdots\circ\theta_{\gamma,\mathfrak d_1}.
\]
注意若 $\gamma$ 穿过若干重叠的射线，排序仍有某种不明确性。然而容易验证：与平行射线相联系的自同构彼此交换，故次序无关紧要。与非平行射线相联系的自同构未必交换，因此需要 $\gamma$ 避开 $\operatorname{Sing}(\mathfrak D)$ 的点。

我们也允许 $\gamma$ 是分片线性的情形，此时 $\gamma'$ 可能在 $t_i$ 处没有定义。在此情形我们要求：若路径 $\gamma$ 与 $\mathfrak d$ 相交，则 $\gamma$ 从 $\mathfrak d$ 的一侧穿到另一侧。此时取 $n_0$ 使得 $\gamma$ 从 $n_0$ 较大的那一侧穿到较小的那一侧。

容易验证 $\theta_{\gamma,\mathfrak D}$ 只依赖于路径 $\gamma$ 在 $M_\R\setminus\operatorname{Sing}(\mathfrak D)$ 中的同伦类。要点是一个简单观察：若 $\gamma$ 是一条连续两次、方向相反地穿过 $\mathfrak D$ 中某条射线的路径，则由两次穿越诱导的自同构互为逆，因而抵消，见图 9。
\end{construction}

% [原文 p.200]
\begin{example}
令 $\mathfrak D=\{(\mathfrak d_1,f_{\mathfrak d_1}),(\mathfrak d_2,f_{\mathfrak d_2}),(\mathfrak d_3,f_{\mathfrak d_3})\}$，其中
\[
  \begin{aligned}
    \mathfrak d_1&=\R r(m_1), & f_{\mathfrak d_1}&=1+c_1w_1z^{m_1},\\
    \mathfrak d_2&=\R r(m_2), & f_{\mathfrak d_2}&=1+c_2w_2z^{m_2},\\
    \mathfrak d_3&=-\R_{\ge0}(r(m_1+m_2)), & f_{\mathfrak d_3}&=1+c_1c_2w_{\mathrm{out}}|r(m_1)\wedge r(m_2)|z^{m_1+m_2},
  \end{aligned}
\]
其中 $m_1,m_2\in T_\Sigma\setminus K_\Sigma$，而 $w_1,w_2$ 与 $w_{\mathrm{out}}$ 分别是 $r(m_1)$、$r(m_2)$ 与 $r(m_1+m_2)$ 的指数。设 $c_1,c_2\in R_k$ 满足 $c_1^2=c_2^2=0$。则当 $\gamma$ 是绕原点的环路时，如图 10，我们可以计算 $\theta_{\gamma,\mathfrak D}$，其中 $\theta_{\gamma,\mathfrak D}=\theta_2^{-1}\theta_3\theta_1^{-1}\theta_2\theta_1$，$\theta_i$ 是来自第一次穿越 $\mathfrak d_i$ 的自同构。注意若用 $\R^2$ 上的标准定向，则对 $i=1,2$，为定义 $\theta_i$ 可取 $n_i=-\frac{1}{w_i}X_{r(m_i)}$：这是零化 $\mathfrak d_i$ 的切空间、且在 $\gamma$ 第一次穿过 $\mathfrak d_1$ 与 $\mathfrak d_2$ 时在 $\gamma'$ 上取负值的余切向量。另一方面，当 $\gamma$ 穿过 $\mathfrak d_3$ 时，$\theta_3$ 用
\[
  n_3=\frac{1}{w_{\mathrm{out}}}X_{r(m_1)+r(m_2)}
\]
定义。于是依次作用 $\theta_1,\theta_2,\theta_1^{-1},\theta_3$ 与 $\theta_2^{-1}$，记住 $c_1^2=c_2^2=0$，我们得到
\[
  \begin{aligned}
    z^m&\mapsto z^m-c_1\langle X_{r(m_1)},r(m)\rangle z^{m+m_1}\\
    &\mapsto z^m-c_2\langle X_{r(m_2)},r(m)\rangle z^{m+m_2}-c_1\langle X_{r(m_1)},r(m)\rangle z^{m+m_1}\\
    &\qquad+c_1c_2\langle X_{r(m_1)},r(m)\rangle\langle X_{r(m_2)},r(m)+r(m_1)\rangle z^{m+m_1+m_2}\\
    &\mapsto z^m-c_2\langle X_{r(m_2)},r(m)\rangle z^{m+m_2}\\
    &\qquad+c_1c_2\bigl(-\langle X_{r(m_2)},r(m)\rangle\langle X_{r(m_1)},r(m)+r(m_2)\rangle\\
    &\qquad\qquad+\langle X_{r(m_1)},r(m)\rangle\langle X_{r(m_2)},r(m)+r(m_1)\rangle\bigr)z^{m+m_1+m_2}\\
    &\mapsto z^m-c_2\langle X_{r(m_2)},r(m)\rangle z^{m+m_2}\\
    &\qquad+c_1c_2\bigl(\langle X_{r(m_1)},r(m)\rangle\langle X_{r(m_2)},r(m_1)\rangle\\
    &\qquad\qquad-\langle X_{r(m_1)},r(m_2)\rangle\langle X_{r(m_2)},r(m)\rangle\\
    &\qquad\qquad+\langle X_{r(m_1)+r(m_2)},r(m)\rangle|r(m_1)\wedge r(m_2)|\bigr)z^{m+m_1+m_2}\\
    &\mapsto z^m+c_1c_2\bigl(\langle X_{r(m_1)},r(m)\rangle\langle X_{r(m_2)},r(m_1)\rangle\\
    &\qquad\qquad-\langle X_{r(m_1)},r(m_2)\rangle\langle X_{r(m_2)},r(m)\rangle\\
    &\qquad\qquad+\langle X_{r(m_1)+r(m_2)},r(m)\rangle|r(m_1)\wedge r(m_2)|\bigr)z^{m+m_1+m_2}.
  \end{aligned}
\]
注意由我们选取的定向，$|r(m_1)\wedge r(m_2)|=\langle X_{r(m_1)},r(m_2)\rangle=-\langle X_{r(m_2)},r(m_1)\rangle$，由此推出上式中 $z^{m+m_1+m_2}$ 的系数为零，于是 $\theta_2^{-1}\theta_3\theta_1^{-1}\theta_2\theta_1$ 是恒等映射。\hfill$\square$
\end{example}

现在由热带树构造散射图，本质上把全部 Maslov 指标零圆盘编码进去。

\begin{definition}
$(X_\Sigma,P_1,\ldots,P_k)$ 中的\textit{热带树}是一条 $d$-标记热带曲线 $h:(\Gamma,p_1,\ldots,p_d)\to M_\R$，满足 $h(p_j)=P_{i_j}$，$1\le i_1<\cdots<i_d\le k$，连同额外数据：一条不加标记的无界边 $E_{\mathrm{out}}\in\Gamma^{[1]}_\infty$ 的选取，使得对任意 $E\in\Gamma^{[1]}_\infty\setminus\{E_{\mathrm{out}}\}$，$h(E)$ 是一个点或是某个 $\rho\in\Sigma^{[1]}$ 的平移。$h$ 的次数 $\Delta(h)$ 定义为不计边 $h(E_{\mathrm{out}})$ 的次数，而 $h(E_{\mathrm{out}})$ 不必是任何 $\rho\in\Sigma^{[1]}$ 的平移。
\end{definition}

% [原文 p.201]
\begin{figure}[htbp]
\centering
\includegraphics[width=0.62\textwidth]{f5_10.png}
\figtitle{10}{}
\end{figure}

$h$ 的 Maslov 指标是
\[
  \mathrm{MI}(h):=2(|\Delta(h)|-d).
\]
给定 $h$ 以及 $E_{\mathrm{out}}$ 内部的一点 $V_{\mathrm{out}}$，我们可以从 $\Gamma$ 中去掉 $E_{\mathrm{out}}\setminus\{V_{\mathrm{out}}\}$ 的无界部分而得到 $\Gamma'$。注意 $V_{\mathrm{out}}$ 是 $\Gamma'$ 的一价顶点。取 $h':\Gamma'\to M_\R$，$h'=h|_{\Gamma'}$。则 $h'$ 是以 $h(V_{\mathrm{out}})$ 为边界、Maslov 指标为 $\mathrm{MI}(h')=\mathrm{MI}(h)$ 的热带圆盘，因为 $|\Delta(h')|=|\Delta(h)|$。

反之，给定热带圆盘 $h':\Gamma'\to M_\R$，可以把 $E_{\mathrm{out}}$ 延拓成无界边，得到热带树 $h:\Gamma\to M_\R$。

如同引理 5.6，标准的热带维数计数论证表明：对一般位置的 $P_1,\ldots,P_k$，热带树 $h$ 在一个维数为 $\mathrm{MI}(h)/2$ 的族中移动。特别地，Maslov 指标零的热带树构成有限集，我们记作 $\operatorname{Trees}(\Sigma,P_1,\ldots,P_k)$。照常，对一般位置的 $P_1,\ldots,P_k$，可以假设所有这些树都是三价的。

\begin{definition}
我们定义 $\mathfrak D(\Sigma,P_1,\ldots,P_k)$ 为这样的散射图：对 $\operatorname{Trees}(\Sigma,P_1,\ldots,P_k)$ 的每个元素 $h$ 含一条射线。对应于 $h$ 的射线形如 $(\mathfrak d,f_{\mathfrak d})$，其中
\begin{itemize}
\item $\mathfrak d=h(E_{\mathrm{out}})$。
\item $f_{\mathfrak d}=1+w_\Gamma(E_{\mathrm{out}})\operatorname{Mult}(h)z^{\Delta(h)}u_{I(h)}$，其中 $u_{I(h)}=\prod_{i\in I(h)}u_i$，而 $I(h)\subseteq\{1,\ldots,k\}$ 由
\[
  I(h):=\{i\mid h(p_j)=P_i \text{ 对某个 } j\}
\]
定义。
\end{itemize}
\end{definition}

\begin{example}
$\mathfrak D(\Sigma,P_1,P_2)$ 如图 11 所示，其中 $\Sigma$ 是 $\PP^2$ 的扇。
\end{example}

\begin{proposition}
设 $P_1,\ldots,P_k$ 取一般位置。若
\[
  P\in\operatorname{Sing}(\mathfrak D(\Sigma,P_1,\ldots,P_k))
\]
是满足 $P\notin\{P_1,\ldots,P_k\}$ 的奇点，且 $\gamma_P$ 是绕 $P$ 的小环路，则
\[
  \theta_{\gamma_P,\mathfrak D(\Sigma,P_1,\ldots,P_k)}=\mathrm{Id}.
\]
\end{proposition}

% [Tutor Note: p.201 命题 5.28 的核心是“奇点 $P$ 处恰有两条射线相撞”的双向对应：(1) 若 $P=\operatorname{Init}(\mathfrak d)$ 且 $\mathfrak d$ 对应树 $h$，在 $V$ 处切开得两个 Maslov 指标非负的圆盘，由 $\mathrm{MI}(h)=\mathrm{MI}(h'_1)+\mathrm{MI}(h'_2)$ 且 $\mathrm{MI}(h)=0$ 得两者均为零；(2) 反向由两棵 Maslov 指标零树在 $P$ 处切开、再粘接并补出新出射边得新树。核对 $I(h_1)\cap I(h_2)=\varnothing$ 的必要性：若相交则两树共享标记点 $P_i$，就无法在 $V$ 处粘接。]

\begin{proof}
设 $P$ 是这样的奇点。假设 $\mathfrak d\in\mathfrak D(\Sigma,P_1,\ldots,P_k)$ 满足 $\operatorname{Init}(\mathfrak d)=P$，并令 $h$ 是对应的树。则 $\Gamma$ 在 $E_{\mathrm{out}}$ 上的唯一顶点 $V$ 满足 $h(V)=P$。注意由于 $P\notin\{P_1,\ldots,P_k\}$，以 $V$ 为顶点的边都不被 $h$ 收缩。此外由一般性假设，$V$ 是三价的，

% [原文 p.202]
\begin{figure}[htbp]
\centering
\includegraphics[width=0.86\textwidth]{f5_11.png}
\figtitle{11}{散射图 $\mathfrak D(\Sigma,P_1,P_2)$。}
\end{figure}

于是若在 $V$ 处切开 $\Gamma$，我们得到两个热带圆盘 $h'_1:\Gamma'_1\to M_\R$ 与 $h'_2:\Gamma'_2\to M_\R$，其边界都是 $P$ 且都有 $V_{\mathrm{out}}=V$。现在 $\mathrm{MI}(h)=\mathrm{MI}(h'_1)+\mathrm{MI}(h'_2)$，所以唯一可能是 $\mathrm{MI}(h'_1)=\mathrm{MI}(h'_2)=0$。于是 $h'_1,h'_2$ 延拓为热带树 $h_i:\Gamma_i\to M_\R$，对应射线 $\mathfrak d_1,\mathfrak d_2$。注意 $P\ne\operatorname{Init}(\mathfrak d_1),\operatorname{Init}(\mathfrak d_2)$ 且 $I(h_1)\cap I(h_2)=\varnothing$。所以满足 $P=\operatorname{Init}(\mathfrak d)$ 的每条射线 $\mathfrak d$ 都来自两条射线 $\mathfrak d_1,\mathfrak d_2$（$P\ne\operatorname{Init}(\mathfrak d_i)$）的相撞。

反之，若给定两条这样的射线 $\mathfrak d_1,\mathfrak d_2$ 经过 $P$，对应树 $h_1,h_2$ 满足 $I(h_1)\cap I(h_2)=\varnothing$，则我们通过下述方式得到一棵新树：在 $P$ 处切开 $h_1$ 与 $h_2$，得到边界在 $P$ 的 Maslov 指标零圆盘 $h'_i:\Gamma'_i\to M_\R$。接着在出射顶点 $V$ 处粘接 $\Gamma'_1$ 与 $\Gamma'_2$，并加上一条以 $V$ 为端点的额外无界边 $E_{\mathrm{out}}$，得到图 $\Gamma$。若 $E_{\mathrm{out},1}$、$E_{\mathrm{out},2}$ 分别是 $\Gamma'_1$、$\Gamma'_2$ 的两条出射边，且 $h'_i(E_{\mathrm{out},i})$ 的、指向远离 $P$ 方向的本原切向量为 $m'_i$，则定义 $h:\Gamma\to M_\R$ 使其在 $\Gamma'_i$ 上限制为 $h'_i$，并把 $E_{\mathrm{out}}$ 映到射线 $P-\R_{\ge0}(w_\Gamma(E_{\mathrm{out},1})m'_1+w_\Gamma(E_{\mathrm{out},2})m'_2)$。取 $w_\Gamma(E_{\mathrm{out}})$ 为 $w_\Gamma(E_{\mathrm{out},1})m'_1+w_\Gamma(E_{\mathrm{out},2})m'_2$ 的指数，则 $h$ 在 $V$ 处平衡。因此 $h$ 是热带树，其 Maslov 指标为零。

作出这些观察之后，为证明命题，定义一个新的散射图 $\mathfrak D_P$，其元素与 $\mathfrak D(\Sigma,P_1,\ldots,P_k)$ 中含 $P$ 的元素一一对应。若 $(\mathfrak d,f_{\mathfrak d})\in\mathfrak D(\Sigma,P_1,\ldots,P_k)$ 是含 $P$ 的射线，则 $\mathfrak D_P$ 中对应元素为 $(\mathfrak d',f_{\mathfrak d})$，其中 $\mathfrak d'$ 在 $P\ne\operatorname{Init}(\mathfrak d)$ 时是 $\mathfrak d$ 的（过原点的）切线，否则是以原点为端点的射线 $\mathfrak d-P$。若 $\gamma_0$ 是与 $\gamma_P$ 定向相同、绕原点的环路，则 $\theta_{\gamma_0,\mathfrak D_P}=\theta_{\gamma_P,\mathfrak D(\Sigma,P_1,\ldots,P_k)}$。

首先考虑最简单的情形：$\mathfrak D_P$ 含两条直线且至多含一条射线。若两条直线对应于树 $h_1$ 与 $h_2$，且 $I(h_1)\cap I(h_2)\ne\varnothing$，则 $h_1$ 与 $h_2$ 不能如上粘接，因为它们经过某个公共标记点 $P_i$。因此 $\mathfrak D_P$ 不含射线。在此情形，与 $\mathfrak d_1$ 和 $\mathfrak d_2$ 相联系的自同构由例 5.24 交换，因为 $u_{I(h_1)}u_{I(h_2)}=0$，于是 $\theta_{\gamma_0,\mathfrak D_P}$ 是恒等映射。

% 边界说明：本块止于 p.202 末句（按 §10 已补完跨页句“与 d_1 和 d_2 相联系的自同构由例 5.24 交换……”。下一块（印刷 p.203 = PDF p.218）应从“若另一方面 $I(h_1)\cap I(h_2)=\varnothing$，……”起续写，不要重复本句。
